% Generated draft: requires mathematical and source review.
\documentclass[11pt]{article}
\usepackage[margin=25mm]{geometry}
\usepackage{fontspec}
\setmainfont{Latin Modern Roman}
\usepackage{amsmath,amssymb,mathtools,unicode-math}
\setmathfont{Latin Modern Math}
\usepackage{enumitem,needspace,xurl,hyperref}
\setlength{\parindent}{0pt}
\setlength{\parskip}{5pt}
\setlength{\emergencystretch}{3em}
\newcommand{\N}{\mathbb{N}}
\newcommand{\Z}{\mathbb{Z}}
\newcommand{\Q}{\mathbb{Q}}
\newcommand{\R}{\mathbb{R}}
\newcommand{\C}{\mathbb{C}}
\newcommand{\F}{\mathbb{F}}
\newcommand{\A}{\mathbb{A}}
\newcommand{\PP}{\mathbb{P}}
\newcommand{\E}{\mathbb{E}}
\newcommand{\eps}{\varepsilon}
\newcommand{\dd}{\,\mathrm{d}}
\begin{document}
\section*{UnsolvedMath ID: 701175}
\textbf{Source problem code:} catalogue entry
\subsection{Definitions and mathematical statement}
A group \(G\) is a \(2\)-group if every element has finite order equal to a power of \(2\). Write
\[
I(G)=\{t\in G:t^2=1,\ t\ne1\}
\]
for its set of involutions. A group \(A\) is periodic if every element has finite order, and locally cyclic if every finitely generated subgroup of \(A\) is cyclic. A group is locally finite if every finitely generated subgroup is finite.

In this draft, a regular group of automorphisms is interpreted in the fixed-point-free sense: for \(A\leq\operatorname{Aut}(G)\),
\[
\{g\in G:a(g)=g\}=\{1\}
\qquad\text{for every }a\in A\setminus\{1\}.
\]
The supplied excerpt does not define ``regular''; this interpretation requires confirmation against the source.

With that qualification, the question attributed to N. M. Suchkov is the following. Let \(G\) be a possibly infinite \(2\)-group. Suppose there exists a periodic, locally cyclic subgroup \(A\leq\operatorname{Aut}(G)\) that is regular in the sense above and acts transitively on \(I(G)\), meaning
\[
\forall t,u\in I(G)\quad\exists a\in A\quad a(t)=u.
\]
Must
\[
\forall\text{ finite }S\subseteq G,\qquad |\langle S\rangle|<\infty
\]
hold? The excerpt does not explicitly require \(I(G)\ne\varnothing\). The displayed transitivity condition is vacuous when \(G\) is trivial; this boundary case already satisfies the conclusion. Every nontrivial \(2\)-group has an involution.
\subsection{Short English statement}
Must a possibly infinite \(2\)-group be locally finite if it admits a regular periodic locally cyclic automorphism group acting transitively on its involutions? Here ``regular'' is provisionally interpreted as requiring every nonidentity automorphism in that group to fix only the identity element of \(G\).
\subsection{Sources}
\begin{itemize}
\item Supplied local UnsolvedMath registry record: permanent ID \(701175\), exact source problem code \texttt{catalogue entry}, title \emph{Kourovka Notebook Problem 15.87}, upstream revision \texttt{d371ee917ac9647fbb06cd6fc529deec6f3dda82}. The record labels the problem ``open''; this is supplied registry metadata, not an independently verified assessment.
\item \emph{The Kourovka Notebook (Unsolved Problems in Group Theory)}, Problem 15.87, proposed by N. M. Suchkov in 2002. The supplied attribution identifies the 15th issue (2002) and the 21st issue edition (2026), \texttt{arXiv:1401.0300v48}. The supplied collection description names E. I. Khukhro and V. D. Mazurov as editors and covers Issues 1--21 (1965--2026). These bibliographic details and the excerpt are supplied source data; the cited edition has not been independently consulted.
\end{itemize}
\end{document}
